\section{The Hypotheses, One by One}\label{sec:recorrido}
This section brings together in a single account what the preceding sections present separately. For each hypothesis, it indicates the observation that motivated it, how it was represented in the model, the results and its status at closure. Table~\ref{tab:recorrido} summarizes the account. A ``compatible'' result means that the model reproduces observations under reasonable assumptions; it does not prove that the event occurred that way. The exploratory compatibility index (ECI) compares hypotheses using the same criteria; it is not a validation.

\subsection{Flow Entry Points Tested in the Model}
\paragraph{H1: Assumed Main-Channel Discharge onto the Alluvial Fan.} \emph{Why it was proposed.} Initial reports collected by Chilevisión and Interferencia placed the origin in the ravine descending to the alluvial fan at Álvaro del Portillo, La Plaza and General Blanche, where the ravine enters a conduit beneath the developed area \citep{chilevision2026,interferencia2026}. The scenario assumes that the conduit receives no flow and assigns the entire inflow to the surface from that point. \emph{How it was represented.} A triangular mixture hydrograph injected at the channel outlet onto the alluvial fan. \emph{Results.} In the model with only the main channel (v01), mud travels along General Blanche and Camino El Alba and reaches close to Padre Hurtado, but it does not reproduce the reported impacts in the southern sector: Av.\ La Plaza Sur, Clínica Universidad de los Andes and the campus remain dry. In the test using held-out observations, it reaches the same 19 of 49 points as the selected hypothesis. \emph{Status.} Inflow through the main channel reproduces the northern branch in these scenarios; by itself it does not reproduce the reported impacts in the southern sector.

\paragraph{H2: Assumed Surface Inflow at the Francisco Bulnes Correa Culvert.} \emph{Why it was proposed.} Interferencia and BioBioChile collected accounts of an overflow where the enclosed channel crosses Francisco Bulnes Correa, at the start of Quebrada Honda, which is among the reported affected streets \citep{interferencia2026,biobio2026origen}. \emph{How it was represented.} A source at the culvert inlet (e03), and hypothetical scenarios of culvert capacity represented as a sink (v10, v11, v17, v18). \emph{Results.} A discharge only at the inlet reaches none of the 49 held-out points: by itself it does not reproduce the full set of observations. With assumed inflow from the alluvial fan, the city model channels about \SI{93000}{m^3} toward the inlet, with peaks near \SI{65}{m^3/s}. In those scenarios, which do not evaluate the real structure, the sink with an assumed capacity of \SI{30}{m^3/s} reaches that limit and simulated flow remains in Quebrada Honda. \emph{Status.} As a sole input, it does not reproduce observations under the assumptions tested. As a transit point for assumed flow from the alluvial fan, it reproduces flow in Quebrada Honda; it does not establish why that street was affected during the event.

\paragraph{H3: Inflow Distributed among Several Channels.} \emph{Why it was proposed.} Sites with reported impacts to the south (La Plaza Sur, the campus, Clínica Universidad de los Andes) are not reached in the model by inflow from the main channel. The mapping used also represents a southern branch reaching the campus and short ravines ending at La Plaza Sur. \emph{How it was represented.} The same hydrograph, distributed in proportion to each catchment's area: \SI{79.5}{\percent} to the main channel, \SI{18.8}{\percent} to the southern branch and \SI{1.7}{\percent} to the short ravines. \emph{Results.} Without hillside rainfall, the configuration with the official southern branch (v21) obtains an ECI of 68/100 and does not reach Clínica Universidad de los Andes. With rainfall on the hillside (H5, below), z02 obtains an ECI of 72/100 and was the hypothesis selected for the held-out test, whose outcome was unfavorable (19 of 49 points). \emph{Status.} It is the basis of all subsequent hypotheses, but by itself leaves much of the official affected area without mud.

\paragraph{H4: Combined Discharge.} \emph{Why it was proposed.} Initial results showed that Quebrada Honda became wet only if part of the simulated flow discharged at the inlet. \emph{How it was represented.} Inflows on the alluvial fan and at the inlet, plus the southern branches (e01 and g03). \emph{Results.} In screening with \SI{150000}{m^3}, Quebrada Honda is wetted at \SIrange{54}{66}{\percent} of its points, but the front remains \SIrange{816}{1007}{m} from Padre Hurtado (Section~\ref{sec:res-calib}). \emph{Status.} It does not meet the reach criterion in those tests.

\paragraph{H5: Rainfall on the Hillside.} \emph{Why it was proposed.} Initial results showed that Clínica Universidad de los Andes remained dry if mud arrived only through the channels. \emph{How it was represented.} H4 plus rainfall on the lower hillside in screening (f02 and f03); subsequently, distributed channels plus rainfall with the full city (v02 and z02). \emph{Results.} z02 obtains an ECI of 72 and reaches 19 of 49 facilities in the unfavorable held-out test. Without the assumed hillside rainfall (v21), the model does not reach Clínica Universidad de los Andes and the ECI falls to 68. \emph{Status.} Assumed rainfall on the lower hillside contributes to reproducing the reported impacts in the southern sector; the held-out test remains unfavorable. Its intensity is revisited below.

\subsection{How Much Mud, and of What Type}
\paragraph{Volume.} \emph{Why it matters.} In the model, reach depends on the assumed volume: with \SI{150000}{m^3}, the front stops \SI{348}{m} from Padre Hurtado; with \SI{300000}{m^3}, it reaches within \SI{100}{m}. \emph{How it was constrained.} Using estimates independent of the urban model (Section~\ref{sec:res-volumen}): a rainfall--runoff balance with soil saturation assumed from the information that about \SI{95}{mm} had fallen at Río Molina between 5 and 7 October \citep{cr2_2026} gives a median of \SI{295000}{m^3}. \emph{Status.} The input volume compatible with the simulated reach agrees in order of magnitude with the rainfall balance under its assumptions; it was not adjusted to make them agree.

\paragraph{Mud Concentration.} \emph{Why it matters.} Thicker mud stops sooner; more fluid mud travels farther and spreads out. \emph{Results.} In the model, thick mud (Glenwood~1) stops \SI{396}{m} from Padre Hurtado; very fluid mud (Aspen~2) obtains 70/100. Quadrature over concentration (Section~\ref{sec:res-posterior}) gives an expected area of \SI{169}{ha} (\SIrange{144.5}{192.3}{ha}) and an ECI between 71 and 73. \emph{Status.} Within the plausible range, concentration has little effect on the comparison between hypotheses.

\subsection{How the City Is Represented in the Model}
\paragraph{Channeling by the Represented Streets.} \emph{Why it was proposed.} The first terrain had 1,103 buildings; mud spread over the surface and did not reproduce the meter reported at Clínica UC Christus \citep{cadena3_2026}. \emph{Results.} With the 42,018 building footprints, mud flows deeper along the streets: next to Clínica UC Christus it rises from \SI{0.3}{m} to \SIrange{1.1}{1.2}{m}, and the culvert inlet receives seven times as much mud. \emph{Status.} Confirmed as a model effect: without the buildings, the reconstruction is not comparable with observations.

\paragraph{Mud Entry into the Represented Houses.} \emph{Why it was proposed.} After the held-out test, the diagnosis showed that the model left the interior of the affected strip dry: mud ran along the streets but did not enter lots, whereas the images consulted are interpreted as showing mud inside houses and courtyards. \emph{How it was represented.} Buildings in the official affected area with floors \SI{0.3}{m} above ground and high resistance, allowing mud to enter and retaining it (v26). \emph{Results.} The fraction of urban mud within the official affected area rises from 0.34 to 0.54, and the verified ECI reaches 76. With half the resistance, the result barely changes: what matters is that mud enters the lots. The tradeoff is the depth next to Clínica UC Christus, which falls to \SI{0.30}{m} because the represented upstream houses store mud in the model. Keeping large blocks solid and making only houses of \SI{250}{m^2} or less porous (d39) does not restore that depth (\SI{0.31}{m}): the effect comes from upstream storage, not building size.
\emph{Status.} It improves agreement with the official affected area, but is a fitting adjustment after the held-out test and has not been validated with new observations.

\subsection{The Universidad de los Andes Campus and La Plaza Sur}
\paragraph{Tested Inflows from the Avenue and Hillside.} \emph{Why it was proposed.} The Meganoticias video is interpreted as showing turbid water descending from the woodland and filling a sunken courtyard on campus (visually estimated depth: \SIrange{15}{35}{cm}) \citep{meganoticias2026uandes}, and El Dínamo reports that Senapred named La Plaza Sur among the most affected streets \citep{eldinamo2026}. The city model left that sector almost dry. \emph{What was tested.} First, assumed surface inflow at Álvaro del Portillo toward the campus: with none of the fluids tested, not even turbid water, does the simulated flow reach the courtyard, because the inflow point lies outside the local catchment delineated for the courtyard. Then, runoff from the hillside itself and from the hill, with features that the \SI{30}{m} terrain does not resolve: the low-lying roadway of La Plaza Sur, the passage from the clearing to the courtyard, and the dirt road on the hill. \emph{Results.} With those represented features, simulated water follows La Plaza Sur northward, enters through the clearing and fills the courtyard to \SIrange{0.12}{0.30}{m}, within the range visually estimated from the video, with assumed rainfall of \SIrange{15}{35}{mm/h}. \emph{Status.} The simulated La Plaza Sur $\rightarrow$ clearing $\rightarrow$ courtyard route is compatible with observations and persists at the rainfall intensities tested. The assumed inflow at Álvaro del Portillo does not reach the courtyard under the assumptions tested. With a fine mesh on the campus within the city simulation and measured rainfall (m36), the courtyard reaches \SI{0.17}{m} and La Plaza Sur is \SI{58}{\percent} wetted: the difference between scales was due to resolution.

\paragraph{More Fluid Flow on the Lower Hillside.} \emph{Why it was proposed.} The cited Meganoticias video is interpreted as showing turbid water, with no discernible urban debris; this motivates testing a more dilute flow on the lower hillside. \emph{Results.} With a lower concentration on the lower hillside, the campus also becomes wet with a fine time step and the verified ECI reaches 78. \emph{Status.} Compatible: the scenario with a more fluid mixture in the southern sector better reproduces the fitting observations; it does not determine the actual flow composition.

\paragraph{The \emph{Bike Park}.} \emph{Why it was proposed.} In the images consulted, the author identified a \emph{bike park} with bare-earth trails above the hill road; its representation was tested to examine the calculated route toward the road. \emph{Results.} With and without the main trail carved, results along the road, on La Plaza Sur, in the clearing and in the courtyard are identical (mrbc and mrbt): with this terrain, water was already draining toward the road. Including the entire area of the \emph{bike park} and upper hillside in the rainfall does increase water input (La Plaza Sur rises from \SI{69}{\percent} to \SI{78}{\percent} wetted). \emph{Status.} In the tests, carving the trail does not change the calculated route; including that area in the applied rainfall increases the model’s water input. Its possible sediment contribution cannot be evaluated with this model, which does not simulate erosion.

\subsection{The Integrated Hypothesis and Rainfall}
\paragraph{Integrated Hypothesis.} It combines the distributed channels, rainfall on the lower hillside, porous houses, more fluid flow on the lower hillside and campus features. With assumed rainfall of \SI{35}{mm/h} for one hour, it obtains 82/100, verified with a fine time step, and meets all four original criteria. With rainfall measured at Río Molina and its recorded temporal distribution, a pulse with a maximum of \SI{10.9}{mm} in \SI{30}{min} (p34b, corrected), it obtains 81/100; the campus reaches \SI{0.49}{m}, La Plaza Sur is wetted at \SI{33}{\percent} of its points and the area with more than \SI{5}{cm} is \SI{222}{ha}. With the same rainfall distributed uniformly over \SI{4}{h} (v33), it falls to 79/100.
\emph{Status.} It is the hypothesis most compatible with the fitting observations, but it was constructed after the held-out test: validating it requires new independent observations.

\paragraph{The Hypothetical Tree Obstacle.} \emph{Why it was proposed.} The author marked a line of trees descending along the main ravine and an accumulation at its end, shortly before the outlet onto the alluvial fan. An assumed obstacle was tested to examine whether the calculated flow is diverted southward. \emph{How it was represented.} A channel carved along the line of trees, with flow injected upstream, and the same channel with the accumulation represented as a \SI{3}{m}-high obstacle (arc and art). \emph{Results.} In the model, the assumed obstacle does not divert the flow southward: the Álvaro del Portillo sector and campus have the same wetted area with and without the assumed obstacle (\SI{20.2}{ha}). The model obstacle retains and spreads mud near the outlet onto the alluvial fan, and the mud travels less far westward. The ECI barely changes (78 versus 79). \emph{Status.} With this terrain, the assumed obstacle at the end of the line of trees does not reproduce the reported impacts in the southern sector.

\paragraph{The Fire.} \emph{Why it was proposed.} According to CR2, the December 2025 fire affected most of the catchment \citep{cr2_2026}. The literature cited in the severity analysis describes possible increases in runoff (hydrophobic soil without cover) and available loose material after fires. \emph{How it was quantified.} Using severity estimated by satellite (dNBR, Section~\ref{sec:res-incendio}): \SI{82}{\percent} of the contributing catchments shows a burn signal, but only \SI{10}{\percent} at moderate-high or high severity. With that severity and the assumed coefficients, the balance for the scenario without fire gives a mud pulse of about \SI{62}{\percent} of the volume with fire (between \SI{51}{\percent} and \SI{76}{\percent}), with \SI{43}{\percent} of the solids. The empirical Staley M1 model gives a median probability of 0.28 with fire versus 0.07 without it, as an illustrative application. \emph{Results in the city.} With the volume without fire (f35b and f38b, corrected), mud wets \SIrange{179}{193}{ha}, compared with \SI{222}{ha} in p34b (\SIrange{13}{19}{\percent} less), reaches fewer streets (0.70--0.73 versus 0.77) and leaves a smaller thickness at Clínica UC Christus (\SIrange{0.19}{0.23}{m} versus \SI{0.30}{m}). However, the ECI does not fall (80--82), because it rewards the fraction of mud within the official affected area, which rises when there is less mud. \emph{Status.} With the assumed coefficients, the scenario with fire calculates a greater volume and, above all, more solids and a higher empirical probability of debris flow. With these metrics, the urban model does not support a claim that this additional volume was necessary to reproduce observations.

\subsection{Aspects Untested at Closure}
The full list appears under ``What Remains Unproven'' (Section~\ref{sec:conclusiones}); three limitations are highlighted here.
\begin{itemize}
\item \textbf{Independent validation.} Everything after the held-out test concerns fitting. Its independent validation remains untested: the integrated hypothesis was not compared with new observations, such as post-event radar images or official high-water marks.
\item \textbf{Clínica UC Christus.} No configuration simultaneously reproduces the meter reported there and the extent within the official affected area.
\item \textbf{What the model does not simulate:} erosion and deposition along the flow path, logs and vehicles swept away, urban drainage, infiltration within the city and changes in flow concentration.
\end{itemize}

\begin{table}[htbp]\centering\small
\caption{Summary of the modeling hypotheses: what motivated them, what was calculated and their status within the tests. ECI: exploratory compatibility index (a comparison of fit, not a measure of predictive accuracy); QH: Quebrada Honda. Compatibility refers to the model; it does not identify causes of the event.}\label{tab:recorrido}
\begin{adjustbox}{max width=\textwidth}\begin{tabular}{@{}>{\raggedright\arraybackslash}p{.17\textwidth}>{\raggedright\arraybackslash}p{.27\textwidth}>{\raggedright\arraybackslash}p{.30\textwidth}>{\raggedright\arraybackslash}p{.22\textwidth}@{}}\toprule
Hypothesis & Motivating information & Model result & Status in the model\\\midrule
H1 main channel & Origin reported on the alluvial fan (press cited in H1) & Reproduces the northern branch; not the southern sector & Insufficient by itself\\
H2 Francisco Bulnes Correa culvert & Overflow reported at the start of Quebrada Honda (press cited in H2) & As sole input, 0 of 49 points; as a flow passage, reproduces Quebrada Honda & Insufficient as sole input\\
H3 distributed inflow & Southern sector outside the main channel & ECI 68 without rainfall (v21); 72 with rainfall (z02, H5) & Basis of subsequent hypotheses\\
H4 combined discharge & Quebrada Honda dry in the model without inflow at the inlet & Screening: QH \SIrange{54}{66}{\percent}; front \SIrange{816}{1007}{m} from Padre Hurtado & Does not meet reach criterion\\
H5 hillside rainfall & Clínica Universidad de los Andes dry in the model without rainfall & ECI 72 (z02); without rainfall, 68; held-out test 19/49 & Compatible fit; unfavorable test\\
Volume & Reach to Padre Hurtado & Assumed \SI{300000}{m^3}, of the order of the balance with saturated soil & Consistent\\
Buildings & Meter reported by Cadena 3 at Clínica UC Christus & Represented streets channel simulated mud & Confirmed in the model\\
Porous houses & Interior of the official affected area dry in the model & ECI 76; Clínica UC Christus falls to \SI{0.30}{m} & Unvalidated fit\\
Campus via La Plaza Sur & Meganoticias video; Senapred according to El Dínamo & Courtyard \SIrange{0.12}{0.30}{m} via the clearing; m36: \SI{0.17}{m} and La Plaza Sur \SI{58}{\percent} & Compatible\\
Overflow from Álvaro del Portillo to campus & Assumed overflow location & Does not reach the courtyard with the fluids tested & Not reproduced in the tests\\
\emph{Bike park} & Bare trails above the road & Carving the trail does not change the simulated route; assumed rainfall on its area adds water & No effect on the simulated route\\
Integrated hypothesis & All of the above & ECI 79--82 depending on rainfall & Best fit, unvalidated\\
Tree obstacle & Line of trees and accumulation identified in images & Does not divert flow southward; retains mud near the alluvial fan & Does not reproduce the southern sector\\
Fire & Fire reported by CR2 nine months earlier & Scenario without fire: about \SI{62}{\percent} of the volume, \SI{43}{\percent} of the solids and about 4-fold lower illustrative model probability; in the city, \SIrange{13}{19}{\percent} less area, ECI 80--82 versus 81 & Sensitivity to assumptions\\\bottomrule
\end{tabular}\end{adjustbox}
\end{table}
